\subsection{A smooth very ample counterexample to full coherent injectivity}

\begin{lemma}[A cyclic cover converts a negative-cohomology class into a restriction kernel]
\label{lem:A-cyclic-coherent-kernel}
\status{Proved; independently audited.}
Let $R$ be a smooth projective geometrically connected $n$-fold over an
algebraically closed field $k$ of characteristic $p$, let $L$ be very
ample, and suppose
$H^{n-1}(R,L^{-1})\ne0$. Fix an embedding given by a space of sections
of $L$. Let $q>1$ be prime to $p$, and choose an ambient homogeneous
form $B$ of degree $q$ whose restriction $b$ has smooth zero divisor
$D$ on $R$. Form the cyclic cover
\[
 \pi:X=\operatorname{Spec}_R\!
       \left(\bigoplus_{j=0}^{q-1}L^{-j}\right)\longrightarrow R,
 \qquad y^q=b,
\]
where the products of terms of total index at least $q$ use
$b:L^{-q}\to\OO_R$. Then $X$ is smooth projective and connected,
$A=\pi^*L$ is very ample, and
\[
 Y=V(y)\subset X,\qquad y\in H^0(X,A),
\]
is a smooth hyperplane section. The map
\[
 H^{n-1}(X,\OO_X)\longrightarrow H^{n-1}(Y,\OO_Y)
\]
is not injective. It also fails to be injective for every smooth
section $V(a y+s)$ with $a\in k^\times$ and $s\in H^0(R,L)$.
\end{lemma}
\begin{proof}
On a trivializing affine open for $L$, the cover has equation
$y^q-b=0$. Away from $D$, its derivative with respect to $y$ is
nonzero. Above $D$, $db$ is nonzero since $D$ is smooth. These
Jacobian computations prove smoothness. The finite morphism $\pi$
is projective. Moreover,
\[
 H^0(X,\OO_X)=\bigoplus_{j=0}^{q-1}H^0(R,L^{-j})=k,
\]
so $X$ is connected; a smooth connected scheme over an algebraically
closed field is integral here, since its irreducible components are
disjoint. The equation $y=0$ identifies $Y$ with $D$.

We check very ampleness rather than infer it from ampleness of a
finite pullback. If $R\subset\mathbf P^N$ is the specified
embedding, the equations of its projective cone in
$\mathbf P^{N+1}$, together with
\[
                 y^q=B(x_0,\ldots,x_N),
\]
define $X$. The cone vertex is excluded by this equation.
On each chart $x_i\ne0$ the scheme is exactly
$\operatorname{Spec}\bigl(\OO_R(R_{x_i})
[y/x_i]/((y/x_i)^q-b/x_i^q)\bigr)$.
Thus this is a closed embedding, and its hyperplane bundle is
$\pi^*L$. The restricted ambient degree-$q$ system separates
first jets, so Bertini supplies the required smooth $D$.
No projective-normality assertion is needed.

Projection formula and finiteness give
\[
 \begin{split}
 H^{n-1}(X,A^{-1})&=\bigoplus_{j=0}^{q-1}
                         H^{n-1}(R,L^{-j-1}),\\
 H^{n-1}(X,\OO_X)&=\bigoplus_{j=0}^{q-1}
                         H^{n-1}(R,L^{-j}).
 \end{split}
\]
Multiplication by $y$ sends the source summand $j=0$ identically
to the target summand $j=1$. Hence it has nonzero image.
Exactness of
$0\to A^{-1}\xrightarrow{\,y\,}\OO_X\to\OO_Y\to0$
puts that image in the kernel in the statement. For $a y+s$ the
same component is $a$ times the identity; the contribution from
$s$ has target index zero and cannot cancel it. Smooth such
sections form a nonempty open, since $y$ itself is smooth.
\end{proof}

\begin{theorem}[Full coherent restriction fails even for smooth degree-one sections]
\label{thm:A-coherent-restriction-false}
\status{Refuted in the unrestricted form; complete counterexamples below.}
For every prime $p$ there is a smooth projective connected
sixfold $X\subset\mathbf P^r_k$, over an algebraically closed
field $k$ of characteristic $p$, and a smooth hyperplane section
$Y$ for which
\[
 \ker\bigl(H^5(X,\OO_X)\longrightarrow H^5(Y,\OO_Y)\bigr)\ne0.
\]
In characteristic two the construction below exhibits a
one-dimensional nonzero subspace of this kernel. Consequently
the unqualified coherent-injectivity assertion in A.1 is false.
This statement does not decide the remaining integral crystalline
claims A.1(a)--(c), or the separate surface case.
\end{theorem}
\begin{proof}
We give all the auxiliary geometry and the required nonvanishing.
Let $V=\mathbf F_p^4$, and let
\[
 Q=\{\,\textstyle\sum_{i=0}^3x_i z_i=0\,\}
       \subset\mathbf P^3_x\times\mathbf P^3_z .
\]
Write $\OO_Q(a,b)$ for the two hyperplane twists. Let
$\mathcal A$ be the kernel of
$V\otimes\OO_{\mathbf P^3_x}\to\OO_{\mathbf P^3_x}(1)$.
For $\alpha:Q=\mathbf P(\mathcal A^\vee)\to\mathbf P^3_x$,
define the rank-two bundle $\mathcal G$ by
\[
 0\longrightarrow\mathcal G\longrightarrow
 \alpha^*\mathcal A^\vee\longrightarrow\OO_Q(0,1)
 \longrightarrow0,
 \qquad \mathcal G'=\mathcal G\otimes\OO_Q(0,1).
\]
Our projective bundles parametrize one-dimensional quotients.
Put
\[
 R=\mathbf P_Q(F_Q^*\mathcal G'),\quad \rho:R\to Q,\quad
 L=\rho^*\OO_Q(1,1)\otimes\OO_\rho(3).
\]
This is a smooth sixfold. If
$\mathcal B=\ker(V^\vee\otimes\OO_Q\to\OO_Q(0,1))$,
the exact sequence
$0\to\OO_Q(-1,0)\to\mathcal B\to\mathcal G\to0$
and the Koszul surjection
$\bigwedge^2V^\vee\otimes\OO_Q(0,-1)\to\mathcal B$
show that $F_Q^*\mathcal G'$ is globally generated by six
sections. Thus $R$ embeds in $Q\times\mathbf P^5$.
The line bundle $L$ is the restriction of $\OO(1,1;3)$
on $\mathbf P^3\times\mathbf P^3\times\mathbf P^5$,
so is very ample.

This is the explicit example of
\citet[Section~2, Claim and proof, arXiv pp.~3--5]{LauritzenRao1997}.
We verify precisely the nonvanishing needed here.
The projective-bundle canonical formula gives
\[
 \omega_R=\rho^*\OO_Q(p-3,p-3)\otimes\OO_\rho(-2).
\]
Serre duality and $\rho_*\OO_\rho(1)=F_Q^*\mathcal G'$
therefore identify
\[
 H^5(R,L^{-1})^\vee
 =H^1\bigl(Q,F_Q^*\mathcal G\otimes
                     \OO_Q(p-2,2p-2)\bigr).
\]
The term $\OO_Q(-2,2p-2)$ is acyclic: use
$0\to\OO(-3,2p-3)\to\OO(-2,2p-2)
\to\OO_Q(-2,2p-2)\to0$ on the product of projective spaces.
The two exact sequences defining $\mathcal B,\mathcal G$,
after Frobenius pullback and twisting, now identify the last
cohomology group with the cokernel of
\[
 \bigoplus_{i=0}^3 H^0(Q,\OO_Q(p-2,2p-2))
 \xrightarrow{(f_i)\mapsto\sum z_i^p f_i}
 H^0(Q,\OO_Q(p-2,3p-2)).
 \tag{A-C1}
\]
Both untwisted terms have no higher cohomology, by the same
incidence sequence. Thus no spectral-sequence assumption is hidden
in this cokernel description.

Here is a direct nonzero functional on that cokernel for every
prime. Set $m=p-2$ and
$T=k[z_0,z_1,z_2,z_3]/(z_0^p,z_1^p,z_2^p,z_3^p)$.
The cokernel is the bidegree $(m,3p-2)$ part of
$k[x_0,\ldots,x_3]\otimes T/(\sum x_i z_i)$.
Pair the $x$-polynomials with auxiliary $u$-polynomials by
apolar differentiation, and pair $T_j$ with $T_{4p-4-j}$
by the coefficient of $z_0^{p-1}\cdots z_3^{p-1}$.
The polynomial
\[
                    (u_0z_1-u_1z_0)^m
\]
defines a nonzero functional: its value on
$x_0^m z_0^{p-1}z_1z_2^{p-1}z_3^{p-1}$ is $m!\ne0$.
It annihilates multiples of $\sum x_i z_i$, since the
transpose of multiplication by this expression is
$\sum z_i\partial/\partial u_i$, which annihilates the
displayed polynomial. It also annihilates the $z_i^p$.
All apolar factorials are invertible because $m<p$.
This proves the required nonzero cokernel.
For $p=2$, (A-C1) is simply the $40$-to-$35$ matrix that
multiplies quadrics by the four squares. Its image consists
of every quartic monomial except $z_0z_1z_2z_3$, so its
cokernel has dimension exactly one.

Apply Lemma~\ref{lem:A-cyclic-coherent-kernel}, with $q=2$
if $p$ is odd and $q=3$ if $p=2$. For complete specification
of the field and branch, embed $R$ by the ambient
$\OO(1,1;3)$ monomials, of which there are
$4\cdot4\cdot\binom83=896$.
Let $k$ be an algebraic closure of
$\mathbf F_p(t_\gamma:|\gamma|=q)$, with one independent
variable for every degree-$q$ monomial in these $896$
coordinates, and set $B=\sum_{|\gamma|=q}t_\gamma Z^\gamma$.
This is the generic ambient form; Bertini proves that
its restriction has smooth zero divisor. The cone equations
and $y^q=B$ now specify $X$ and $Y=(y=0)$ completely.
The preceding lemma supplies the asserted nonzero kernel.

The retained program \path{code/coherent_restriction_counterexample.py}
checks the apolar functional against every relevant monomial
incidence relation for $p=2,3,5,7,11,13$, checks the exact
characteristic-two rank $34$, and checks the cyclic
multiplication identity. Its output states its scope:
the all-prime proof and generic smoothness are proved above,
not inferred from finite samples.
\end{proof}

\begin{remark}[Degrees and the surviving integral question]
\label{rem:A-cyclic-degree-scope}
\status{Proved under stated hypotheses.}
In this fixed embedding the failure occurs at degree one,
on a nonempty open set of smooth sections. For every
$d\geq\max(1,\operatorname{reg}\omega_X)$ coherent restriction
is injective by Theorem~\ref{thm:A-high-degree}.
At any specified intermediate degree its kernel is exactly
the image of the finite-dimensional multiplication map
\[
 H^5(X,A^{-d})\xrightarrow{\,s\,}H^5(X,\OO_X),
\]
which can be computed from the displayed finite cyclic algebra
and the two projective-bundle sequences. The exact intermediate
list for this sixfold is not claimed. The surface question is
settled separately by Theorem~\ref{thm:OC-surface-coherent-counterexample}
below. The crystalline torsion and saturation obligations remain
open outside the proved ranges.
\end{remark}
